\documentclass[12pt]{article}
\usepackage[T1]{fontenc}
\usepackage[utf8]{inputenc}
\usepackage{lmodern}
\usepackage{textcomp}
\usepackage{enumitem}
\usepackage{geometry}
\geometry{margin=1in}
\begin{document}
\begin{center}
Answer Key - Psychology - Undergraduate Level Assessment 7 \
\small (Answers are ordered exactly as in the question paper.)
\end{center}

\section*{Multiple Choice}
\begin{enumerate}[label=\arabic*.]
\item \textbf{Answer:} A
\\ \textit{Options:} A) lead to a return in the opposite direction; B) reduce the attractiveness of the outcome approached; C) lead to a compromise between the two outcomes; D) magnify the conflict in that direction; E) create a new conflict in the opposite direction
\\ \textit{Note:} In an approach-approach conflict, individuals face a choice between two desirable options, and once they commit to one option, they may experience a psychological pull towards the other option, often resulting in a return or reconsideration towards the alternative choice.
\item \textbf{Answer:} A
\\ \textit{Options:} A) For reducing activity levels a higher dose of the stimulant is required.; B) Stimulant therapy is only effective for adult ADHD, not for children.; C) A higher dose of the treatment leads to severe physical dependency.; D) A limitation of stimulant therapy is that not all children improve.; E) The treatment does not improve attention span.
\\ \textit{Note:} The statement is incorrect because higher doses of CNS stimulants may not necessarily lead to a reduction in activity levels and can instead lead to increased side effects and diminished effectiveness.
\item \textbf{Answer:} A
\\ \textit{Options:} A) None of the above; B) II and III only; C) I and III only; D) I, II and III; E) III only
\\ \textit{Note:} The correct answer is "None of the above" because Harlow's study, while providing insights into attachment and maternal care, specifically focused on non-human primates and does not directly translate to human behavior or development due to significant biological and social differences between species.
\item \textbf{Answer:} A
\\ \textit{Options:} A) behavior therapy; B) Exposure therapy; C) cognitive therapy; D) psychoanalysis; E) Rational emotive behavior therapy
\\ \textit{Note:} The correct answer is behavior therapy because the therapist's affirmative nodding and verbal cues indicate active listening and reinforcement of Amy's behavior in expressing her thoughts and feelings, which are key components of behavior therapy techniques.
\item \textbf{Answer:} C
\\ \textit{Options:} A) adjustment; B) homeostasis; C) exhaustion; D) stressor detection; E) resolution
\\ \textit{Note:} The exhaustion stage of Selye's General Adaptation Syndrome (GAS) theory is when the body's resources are depleted after prolonged stress, leading to increased vulnerability to disease and illness.
\end{enumerate}

\section*{True / False}
\begin{enumerate}[label=\arabic*.]
\setcounter{enumi}{5}
\item \textbf{Answer:} False
\\ \textit{Options:} A) True; B) False
\\ \textit{Note:} The statement is false because color blindness is primarily a recessive trait linked to the X chromosome, which means a woman must have two copies of the color blindness gene to transmit it, while her son only needs one copy from her to express the trait.
\item \textbf{Answer:} False
\\ \textit{Options:} A) True; B) False
\\ \textit{Note:} The statement is false because while nonparental early care and education can influence language development, research indicates that the quality of parent-child interactions and the home environment plays a more critical role in the development of vocal language skills.
\item \textbf{Answer:} False
\\ \textit{Options:} A) True; B) False
\\ \textit{Note:} The statement is false because having accidents often results from risky behaviors or environments, which can be associated with lower physical fitness rather than higher fitness levels among workers.
\item \textbf{Answer:} True
\\ \textit{Options:} A) True; B) False
\\ \textit{Note:} The Thematic Apperception Test, developed in the 1930 s by Henry Murray and colleagues, is considered one of the first standardized measures of psychological assessment because it systematically evaluates individuals' projections of their own thoughts and feelings through storytelling based on ambiguous images.
\item \textbf{Answer:} True
\\ \textit{Options:} A) True; B) False
\\ \textit{Note:} Kohler's studies with apes illustrated that these primates could solve complex problems by suddenly understanding the solution rather than through a gradual process of trial and error, thereby demonstrating the concept of insight learning.
\end{enumerate}

\section*{Long Form Response}
\begin{enumerate}[label=\arabic*.]
\setcounter{enumi}{10}
\item \textbf{Answer:} Shadowing is a cognitive technique employed in attention studies where a participant listens to a series of words and is instructed to repeat them aloud immediately. This method is significant as it allows researchers to investigate the processing of auditory information and the mechanisms of selective attention. By requiring subjects to focus on one auditory channel while ignoring others, shadowing provides insights into how individuals manage competing stimuli and allocate cognitive resources. The implications of this technique are profound, enhancing our understanding of cognitive processes such as memory, perception, and the dynamics of attention in complex environments. Ultimately, shadowing serves as a crucial tool in elucidating the intricacies of human cognition and the underlying neural mechanisms that facilitate attention.
\\ \textit{Note:} correct because it accurately describes shadowing as a technique that enables researchers to study selective attention and auditory processing by requiring participants to focus on specific stimuli while disregarding others, thus illuminating cognitive resource allocation.
\item \textbf{Answer:} The DSM-5 outlines specific diagnostic criteria for problem gambling, which includes a pattern of gambling behavior characterized by persistent and recurrent problematic gambling leading to significant impairment or distress. A critical component of this diagnosis is the concept of loss of control, which refers to the individual's inability to limit the amount or frequency of gambling despite repeated attempts to do so. This loss of control is significant as it highlights the compulsive nature of the behavior, distinguishing problem gambling from occasional or recreational gambling. Furthermore, the inability to regulate gambling can lead to various negative consequences, including financial hardships, relationship issues, and emotional distress, emphasizing the need for effective intervention and treatment. Understanding loss of control is essential in framing problem gambling as a behavioral addiction, necessitating a comprehensive approach to address its psychological and social implications.
\\ \textit{Note:} correct because it accurately identifies the DSM-5 criteria for problem gambling, emphasizing the role of loss of control as a central feature that distinguishes this disorder and underscores its compulsive nature.
\end{enumerate}

\vfill
\noindent\textit{End of Answer Key}
\end{document}